\section{A half-dimensional sample with a common improving family}

The following statements retain the sample, its common operator image, and
its collective trim. They concern one fixed input space, rather than an
adaptive sequence. Their role in the resampling description of
\cite[lines 749--755]{Arad2016RigorousRG}, and the distinction from the
sample-once Figure~\texttt{fig:merge-prime}, are explained
in~\cite{gap:alvv16_bootstrap_constants}.

\begin{lemma}[Numerical margins for sampling and trimming]
    \label{lem:half_sampling_margins}
    \lean{Real.half_sampling_filter_error_le, Real.half_sampling_transfer_error_le, Real.half_sampling_trim_error_le}
    \leanok
    If $d\ge2$, $d^{12}\Delta\le10^{-5}$, and
    $\alpha\ge0.03/d^2$, then $\Delta/\alpha^2\le10^{-4}$.
    If $\alpha>0$ and $\varepsilon\le\alpha^2/100$, then
    \[
      (1+\alpha/4)(1-\alpha)+(1+4/\alpha)\varepsilon
        \le1-\alpha/2.
    \]
    Finally, if $\delta\le27/2500$, $\varepsilon\le10^{-5}$, and
    $z\le10^{-3}$, then
    \[
      (1+1/20)\bigl((1+1/20)\delta+21\varepsilon+2z\bigr)
        +21\varepsilon\le3/200.
    \]
\end{lemma}
\begin{proof}\leanok
    The first inequality follows from $d^8\ge256$ and the lower bound on
    $\alpha$. The other two follow by expansion.
\end{proof}

\begin{lemma}[Unconditional dimension decrease]
    \label{lem:half_sampling_image_dimension}
    \lean{Submodule.finrank_half_sample_image_le}
    \leanok
    Let $W\subseteq M$ have dimension $q$, fix an orthonormal basis of $W$,
    and let $(A_i)_{i\in I}$ be operators on $M$ with $|I|\le D^2$.
    Put $s=\lfloor q/(2D^2)\rfloor$, with $s=0$ when $D=0$.
    Every unitary rotation of the first $s$ basis vectors spans a space
    $S\subseteq W$ of dimension $s$, and
    $\dim\operatorname{span}_{i\in I}A_iS\le q/2$.
    No target or viability assumption is required.
\end{lemma}
\begin{proof}\leanok
    The image dimension is at most $|I|s\le D^2s\le q/2$.
\end{proof}

\begin{theorem}[One original-target event for every improving target]
    \label{thm:half_sampling_uniform_image_event}
    \lean{Submodule.exists_measurable_half_sample_uniform_image_event}
    \leanok
    \uses{lem:half_sampling_margins, lem:half_sampling_image_dimension}
    Let $T_0,T_x\subseteq M\otimes F$, where the parameter set is nonempty,
    and let $D\ge2$. Suppose $T_0$ is $\varepsilon$-close to every $T_x$,
    every $T_x$ is $1/200$-close to $T_0$, and
    $\varepsilon\le(0.03/D^2)^2/100$. A single family $(A_i)_{i\in I}$,
    $|I|\le D^2$, improves every $\delta$-viable space for every $T_x$ to
    error $\Delta/(1-\delta)^2$, for all $\delta<1$.
    Assume $D^{12}\Delta\le10^{-5}$.

    Let $W$ be $0.04$-viable for $T_0$, with dimension $q\ge4D^2$ and a
    specified orthonormal basis. Put
    $s=\lfloor q/(2D^2)\rfloor$ and $\alpha=0.96s/(8q)$.
    Write $S_U$ for the literal Haar sample and
    $Y_U=\operatorname{span}_{i\in I}A_iS_U$.
    There is a Borel event $E$ whose complement has Haar measure at most
    \[
       \eta=(1+2/\sqrt\alpha)^{2\dim T_0}q\,e^{-s/16},
    \]
    and whose measure is at least $1-\eta$, such that for every $U\in E$,
    $S_U$ is $(1-\alpha)$-viable for $T_0$ and
    $(1-\alpha/2)$-viable for every $T_x$; the same $Y_U$ is
    $4\Delta/\alpha^2$-viable for every $T_x$ and $27/2500$-viable for $T_0$.
    For every $U$, including those outside $E$, $\dim Y_U\le q/2$.
    The outer measure of failure of the last original-target conclusion
    is also at most $\eta$.
\end{theorem}
\begin{proof}\leanok
    Sample against the fixed target $T_0$. Weighted composition with
    parameter $\alpha/4$ transfers this same sample to all $T_x$.
    Apply the common family, then transfer back through one fixed $T_x$.
    The numerical margins and the unconditional dimension estimate give
    the assertions. No measurability of an intersection over the parameters
    is asserted or needed.
\end{proof}

\begin{theorem}[Reference trimming of the same sampled image]
    \label{thm:half_sampling_original_reference_trim}
    \lean{Entropy.exists_measurable_half_sample_reference_trim_event}
    \leanok
    \uses{thm:half_sampling_uniform_image_event, lem:half_sampling_margins}
    In the preceding setting, use coordinate tensor factors and an increasing
    list of $m$ cuts. Suppose a reference space $R$ is mutually
    $\varepsilon_R$-close to $T_0$, with $0\le\varepsilon_R\le10^{-5}$,
    and every reference vector has Schmidt rank at most $b$ at each cut.
    Choose $\xi>0$ with $\xi\sqrt{mb}\le10^{-3}$.
    Collectively trim the specified orthonormal basis of the same $Y_U$
    at threshold $\xi$ to obtain $V_U$.
    On the same event $E$, $V_U$ is $3/200$-viable for $T_0$.
    For every $U$, $\dim V_U\le\dim Y_U\le q/2$.
    The outer measure of original-target failure is at most $\eta$.
    All preceding conclusions about $Y_U$ remain in force.
\end{theorem}
\begin{proof}\leanok
    Apply weighted reference trimming with parameter $1/20$ and then
    Lemma~\ref{lem:half_sampling_margins}. Trimming is a linear image of
    the space being trimmed, so its dimension cannot increase.
\end{proof}

\begin{theorem}[The canonical sampling step under product Haar measure]
    \label{thm:rank_product_half_sample_reference_trim}
    \lean{Entropy.exists_measurable_randomRangeSubspace_half_trim_event}
    \leanok
    \uses{def:rank_unitary_product_haar,thm:half_sampling_original_reference_trim}
    Retain the targets, common family, cuts and reference hypotheses of
    Theorem~\ref{thm:half_sampling_original_reference_trim}. Let
    $I:\mathbb C^d\to M$ be induced by a finite coordinate equivalence.
    For any matrix $P\in M_d(\mathbb C)$, put
    $q=\operatorname{rank}P$ and $W=I(\operatorname{range}P)$.
    Assume $q\ge4D^2$ and that $W$ is $0.04$-viable for $T_0$.

    For $U\in\Omega_d$, use the coordinate $U_q$ at the actual rank to sample
    $s=\lfloor q/(2D^2)\rfloor$ columns in the canonical orthonormal range
    basis of $P$, and transport this literal sample by $I$. Denote it by
    $S_U$, put $Y_U=\operatorname{span}_iA_iS_U$, and let $V_U$ be the
    same collective reference trim of $Y_U$. Set
    \begin{align}
        \alpha&=\frac{0.96s}{8q},&
        \eta&=(1+2/\sqrt\alpha)^{2\dim T_0}q\,e^{-s/16}.
    \end{align}
    Then $s\ge1$ and $0<\alpha\le1$. Every seed satisfies
    $S_U\subseteq W$ and $\dim S_U=s$. Both $Y_U$ and $V_U$ have
    dimension at most $q/2$.

    There is a Borel event $E\subseteq\Omega_d$ with
    $\mu_d(E^c)\le\eta$ and $\mu_d(E)\ge1-\eta$ such that, on $E$,
    $Y_U$ is $4\Delta/\alpha^2$-viable for every $T_x$ and
    $27/2500$-viable for $T_0$, while $V_U$ is $3/200$-viable for $T_0$.
    The outer measure of failure of this last conclusion is at most $\eta$.
\end{theorem}
\begin{proof}\leanok
    \uses{thm:half_sampling_original_reference_trim,thm:rank_unitary_haar_marginal}
    Map the canonical range basis through the restriction of $I$ from
    $\operatorname{range}P$ onto $W$. Sampling in this mapped basis is
    exactly the image under $I$ of the original canonical sample.
    Apply Theorem~\ref{thm:half_sampling_original_reference_trim} once
    with this basis, and take the inverse image of its one Borel event
    under $U\mapsto U_q$. The coordinate Haar marginal preserves the
    measures of both the event and its complement. The sample, common
    image and trim agree pointwise with the stated operations, so all
    dimension and viability conclusions transfer unchanged. The true
    output-failure set is contained in $E^c$.
\end{proof}

The law $\mu_d$ is fixed before $P$. The event is obtained at that fixed
input; no measurable choice of the events as $P$ varies, or joint
probability estimate for an adaptive sequence, is asserted.
